\chapter{Thực nghiệm và đánh giá}

\section{Mục tiêu thực nghiệm}

Chương này trình bày kế hoạch thực nghiệm và khung đánh giá cho hệ thống KGsAuto. Trọng tâm đánh giá là ứng dụng Knowledge Graph trong hệ thống hỏi đáp, thông qua việc so sánh các chế độ truy xuất Semantic Search, Naive Graph RAG, Graph Search và Hybrid RAG bằng khung RAGAS.

Đối với các giai đoạn trích xuất Knowledge Graph và Entity Resolution, khóa luận chỉ báo cáo thống kê vận hành do chưa có tập ground truth phù hợp để đánh giá định lượng độ chính xác. Các kết quả chi tiết sẽ được bổ sung sau khi quá trình chạy và kiểm tra thực nghiệm hoàn tất.

\section{Dữ liệu và thiết lập thực nghiệm}

Dữ liệu thực nghiệm là tập tài liệu Markdown về Trường Đại học Công nghệ. Tập dữ liệu này được sử dụng cho hai mục đích: xây dựng Knowledge Graph thông qua pipeline trích xuất và Entity Resolution, đồng thời tạo tập câu hỏi đánh giá cho hệ thống hỏi đáp.

Bảng~\ref{tab:experiment_setup_placeholder} trình bày các thông tin thiết lập thực nghiệm dự kiến. Các giá trị cụ thể có thể được cập nhật sau khi chốt lần chạy cuối cùng.

\begin{table}[H]
\centering
\caption{Thiết lập thực nghiệm dự kiến}
\label{tab:experiment_setup_placeholder}
\begin{tabularx}{\textwidth}{|p{0.36\textwidth}|X|}
\hline
\textbf{Thành phần} & \textbf{Giá trị / ghi chú} \\
\hline
Nguồn dữ liệu & Tài liệu Markdown về Trường Đại học Công nghệ \\
\hline
Graph database & Neo4j \\
\hline
Vector database & Qdrant \\
\hline
Embedding model & Sẽ cập nhật sau khi chốt cấu hình \\
\hline
LLM cho RAG runtime & Sẽ cập nhật sau khi chốt cấu hình \\
\hline
LLM/RAGAS evaluator & Sẽ cập nhật sau khi chốt cấu hình \\
\hline
Tập câu hỏi đánh giá & Sẽ cập nhật sau khi sinh/chọn tập đánh giá cuối cùng \\
\hline
\end{tabularx}
\end{table}

\section{Thống kê pipeline xây dựng Knowledge Graph}

Phần này sẽ báo cáo các thống kê vận hành của pipeline trích xuất và Entity Resolution. Các thống kê này nhằm mô tả quy mô xử lý và artifact sinh ra, không được diễn giải như độ chính xác của extraction hoặc Entity Resolution.

\begin{table}[H]
\centering
\caption{Thống kê dự kiến cho giai đoạn trích xuất Knowledge Graph}
\label{tab:extraction_stats_placeholder}
\begin{tabularx}{\textwidth}{|p{0.45\textwidth}|X|}
\hline
\textbf{Chỉ số} & \textbf{Giá trị} \\
\hline
Số tài liệu Markdown đầu vào & TBD \\
\hline
Số tài liệu xử lý thành công & TBD \\
\hline
Số node được trích xuất & TBD \\
\hline
Số relationship được trích xuất & TBD \\
\hline
Số node trung bình mỗi tài liệu & TBD \\
\hline
Số relationship trung bình mỗi tài liệu & TBD \\
\hline
Thời gian xử lý / chi phí token & TBD \\
\hline
\end{tabularx}
\end{table}

\begin{table}[H]
\centering
\caption{Thống kê dự kiến cho giai đoạn Entity Resolution}
\label{tab:er_stats_placeholder}
\begin{tabularx}{\textwidth}{|p{0.45\textwidth}|X|}
\hline
\textbf{Chỉ số} & \textbf{Giá trị} \\
\hline
Số thực thể trước Entity Resolution & TBD \\
\hline
Số cụm ứng viên & TBD \\
\hline
Số noise item & TBD \\
\hline
Silhouette score & TBD \\
\hline
Số canonical entity sau Entity Resolution & TBD \\
\hline
Số relationship trước rewiring & TBD \\
\hline
Số relationship sau rewiring & TBD \\
\hline
Số relationship trùng lặp được loại bỏ & TBD \\
\hline
\end{tabularx}
\end{table}

\section{Đánh giá khả năng xác minh của triplet trong Knowledge Graph}

Bên cạnh các thống kê vận hành của pipeline, khóa luận thực hiện một bước đánh giá định tính có cấu trúc đối với các quan hệ trong Knowledge Graph sau khi được nhập vào Neo4j. Mục tiêu của bước này là ước lượng mức độ các triplet trong graph có thể được xác minh từ thông tin hiện có, thay vì chỉ báo cáo số lượng node và relationship được sinh ra.

Quy trình đánh giá được thực hiện như sau. Trước hết, hệ thống lấy ngẫu nhiên 100 triplet có dạng \texttt{(subject)-[predicate]->(object)} từ Neo4j. Với mỗi triplet, hệ thống xây dựng một đối tượng đánh giá gồm tên, nhãn và thuộc tính của subject, predicate, object, đồng thời truy xuất các trường bằng chứng cục bộ trong graph như \texttt{description}, \texttt{summary}, \texttt{text}, \texttt{source}, \texttt{chunk\_id} hoặc \texttt{doc\_id} nếu có. Đối tượng này sau đó được đưa cho LLM judge để phân loại vào một trong ba nhãn: \texttt{supported}, \texttt{contradicted} và \texttt{not\_enough\_info}.

Ba nhãn được hiểu như sau. Nhãn \texttt{supported} cho biết thông tin trong graph hỗ trợ triplet được đánh giá. Nhãn \texttt{contradicted} cho biết bằng chứng hiện có mâu thuẫn với triplet. Nhãn \texttt{not\_enough\_info} được dùng khi thông tin trong graph còn thiếu, mơ hồ hoặc không đủ mạnh để xác minh triplet. Cách phân loại này giúp tách biệt hai loại vấn đề khác nhau: quan hệ sai rõ ràng và quan hệ chưa đủ bằng chứng để xác nhận.

Bảng~\ref{tab:triplet_verifiability_results} trình bày kết quả đánh giá trên 100 triplet được lấy ngẫu nhiên. Trong lần chạy này, toàn bộ 100 phản hồi của LLM judge đều hợp lệ, không có phản hồi lỗi định dạng.

\begin{table}[H]
\centering
\caption{Kết quả đánh giá khả năng xác minh của 100 triplet ngẫu nhiên trong Knowledge Graph}
\label{tab:triplet_verifiability_results}
\begin{tabularx}{\textwidth}{|p{0.32\textwidth}|p{0.18\textwidth}|p{0.18\textwidth}|X|}
\hline
\textbf{Nhãn đánh giá} & \textbf{Số lượng} & \textbf{Tỷ lệ} & \textbf{Diễn giải} \\
\hline
\texttt{supported} & 86 & 86\% & Triplet được bằng chứng trong graph hỗ trợ \\
\hline
\texttt{contradicted} & 8 & 8\% & Triplet có dấu hiệu mâu thuẫn với bằng chứng hiện có \\
\hline
\texttt{not\_enough\_info} & 6 & 6\% & Graph chưa có đủ bằng chứng để xác minh triplet \\
\hline
\end{tabularx}
\end{table}

Kết quả cho thấy phần lớn triplet trong mẫu ngẫu nhiên có thể được xác minh từ thông tin hiện có trong graph, với \textit{micro supported rate} đạt 86\%. Khi tính trung bình theo từng loại predicate, \textit{macro supported rate} đạt 84.74\%, cho thấy kết quả không chỉ đến từ một loại quan hệ chiếm ưu thế trong mẫu. Tuy nhiên, 8\% triplet bị đánh giá là mâu thuẫn và 6\% triplet chưa đủ thông tin cho thấy graph vẫn còn lỗi hoặc thiếu provenance ở một số quan hệ. Các trường hợp này là tín hiệu quan trọng để phân tích lỗi ở pipeline extraction, Entity Resolution hoặc bước import vào Neo4j.

Cần lưu ý rằng kết quả trên không thay thế cho một benchmark có ground truth thủ công. Đây là đánh giá dựa trên LLM judge và bằng chứng cục bộ trong graph, do đó phù hợp để quan sát xu hướng chất lượng và phát hiện nhóm lỗi cần kiểm tra sâu hơn. Trong các lần đánh giá tiếp theo, có thể mở rộng phương pháp bằng cách lấy mẫu phân tầng theo predicate, bổ sung kiểm tra thủ công một phần mẫu và truy xuất thêm provenance từ tài liệu gốc.

\section{Đánh giá hệ thống hỏi đáp bằng RAGAS}

Phần đánh giá chính tập trung vào hệ thống hỏi đáp. Bốn chế độ truy xuất được so sánh gồm Semantic Search, Naive Graph RAG, Graph Search và Hybrid RAG. Các metric dự kiến gồm faithfulness, context precision, context recall và latency.

Bảng~\ref{tab:ragas_results_placeholder} là bảng kết quả dự kiến để điền sau khi hoàn tất đánh giá bằng RAGAS.

\begin{table}[H]
\centering
\caption{Bảng kết quả đánh giá hỏi đáp bằng RAGAS}
\label{tab:ragas_results_placeholder}
\begin{tabularx}{\textwidth}{|p{0.22\textwidth}|p{0.14\textwidth}|p{0.16\textwidth}|p{0.18\textwidth}|X|}
\hline
\textbf{Chế độ} & \textbf{Latency TB} & \textbf{Faithfulness} & \textbf{Context precision} & \textbf{Context recall} \\
\hline
Semantic Search & TBD & TBD & TBD & TBD \\
\hline
Naive Graph RAG & TBD & TBD & TBD & TBD \\
\hline
Graph Search & TBD & TBD & TBD & TBD \\
\hline
Hybrid RAG & TBD & TBD & TBD & TBD \\
\hline
\end{tabularx}
\end{table}

Ngoài bảng tổng hợp, phần đánh giá có thể bổ sung phân tích theo nhóm câu hỏi, ví dụ câu hỏi single-hop, multi-hop, câu hỏi về đơn vị, con người, chương trình đào tạo hoặc sự kiện. Nếu có đủ dữ liệu, phần phân tích lỗi sẽ chỉ ra các trường hợp context recall thấp, context precision thấp hoặc câu trả lời chưa bám sát bằng chứng.

\section{Phân tích kết quả và thảo luận}

Phần này sẽ được hoàn thiện sau khi có kết quả đánh giá cuối cùng. Các nội dung thảo luận dự kiến gồm: so sánh vai trò của văn bản và graph trong các chế độ RAG, trade-off giữa chất lượng và độ trễ, ảnh hưởng của chất lượng graph đến Graph Search và Hybrid RAG, cũng như các lỗi thường gặp trong quá trình truy xuất bằng chứng.

\section{Tóm tắt chương}

Chương này hiện đóng vai trò là khung thực nghiệm và đánh giá. Các bảng kết quả sẽ được cập nhật sau khi quá trình chạy thực nghiệm và chấm điểm bằng RAGAS hoàn tất.
